\subsection{两个模的张量积}

设 \(G_{1}, G_{2}\) 为两个 Z-模；从集合 \(G = G_{1} \times G_{2}\) 到某个 Z-模的映射 u 称为双加性的（或 Z-双线性的），如果 \(u(x_{1}, x_{2})\) 关于 \(x_{1}\) 和关于 \(x_{2}\) 都是“加性的”；确切地说，这意味着，对于 \(x_{1}, y_{1}\) 在 \(G_{1}, x_{2}, y_{2}\) 在 \(G_{2}\) ，

\[
\begin{array}{l} u (x _ {1} + y _ {1}, x _ {2}) = u (x _ {1}, x _ {2}) + u (y _ {1}, x _ {2}) \\ u (x _ {1}, x _ {2} + y _ {2}) = u (x _ {1}, x _ {2}) + u (x _ {1}, y _ {2}). \end{array}
\]

注意这特别蕴含 \(u(0, x_2) = u(x_1, 0) = 0\) 对所有 \(x_1 \in \mathbf{G}_1, x_2 \in \mathbf{G}_2\) 成立。

设 A 为一个环，E 为一个右 A-模，F 为一个左 A-模。我们将考虑如下泛映射问题（集合论，IV，§ 3，第 1 号），其中 \(\Sigma\) 为 Z-模结构的类（此时态射即为 Z-线性映射，换言之，加法群同态），而 \(\alpha\) -映射是映射 \(f\) （从 E \(\times\) F 到某个 Z-模 G ），它们是 Z-双线性的，并且进一步满足：对所有 \(x \in \mathbf{E}\) , \(y \in \mathbf{F}\) 和 \(\lambda \in \mathbf{A}\)

\[
f (x \lambda , y) = f (x, \lambda y).\tag{1}
\]

我们证明该问题有解。为此，我们考虑Z-模 \(C = \mathbf{Z}^{(E \times F)}\) ，即由 \(E \times F\) 的元素以Z中为系数的形式线性组合所构成的Z-模（§1，第11号），其一组基可视为由有序对 \((x, y)\) （其中 \(x \in E\) ， \(y \in F\) ）组成。设D为C的子Z-模，由以下类型的元素生成：

\[
\left\{ \begin{array}{l} (x _ {1} + x _ {2}, y) - (x _ {1}, y) - (x _ {2}, y) \\ (x, y _ {1} + y _ {2}) - (x, y _ {1}) - (x, y _ {2}) \\ (x \lambda , y) - (x, \lambda y) \end{array} \right.\tag{2}
\]

其中 \(x, x_{1}, x_{2}\) 属于E， \(y, y_{1}, y_{2}\) 属于F，且 \(\lambda\) 属于A。

定义1. 右A-模E与左A-模F的张量积，记为 \(\mathbf{E} \otimes_{\mathbf{A}} \mathbf{F}\) 或 \(\mathbf{E} \otimes_{\mathbf{A}} \mathbf{F}\) （若不致引起混淆，也简记为 \(\mathbf{E} \otimes \mathbf{F}\) ），指的是商 \(\mathbf{Z}\) -模 C/D（即 \(\mathbf{Z}\) -模 C——由 \(\mathbf{E} \times \mathbf{F}\) 的元素以 \(\mathbf{Z}\) 中为系数的形式线性组合构成——除以由 (2) 中某一类型的元素所生成的子模 \(\mathbf{D}\) 而得的商）。对于 \(x \in \mathbf{E}\) 和 \(y \in \mathbf{F}\) ， \(\mathbf{E} \otimes_{\mathbf{A}} \mathbf{F}\) 中作为元素 \((x, y)\) （属于 \(\mathbf{C} = \mathbf{Z}^{(\mathbf{E} \times \mathbf{F})}\) ）的典范像的那个元素记为 \(x \otimes y\) ，并称为 \(x\) 与 \(y\) 的张量积。

映射 \((x,y)\mapsto x\otimes y\) 于 \(\mathbf{E}\times \mathbf{F}\) 到 \(\mathbf{E}\otimes_{\mathbf{A}}\mathbf{F}\) 称为典范的。它是一个满足条件(1)的Z-双线性映射。

我们证明张量积 \(E \otimes_{A} F\) 并且上述典范映射构成了先前提出的泛映射问题的一个解。确切地说：

命题1. (a) 设 \(g\) 是一个 \(\mathbf{Z}\) -线性映射（从 \(\mathbf{E} \otimes_{\mathbf{A}} \mathbf{F}\) 到某个 \(\mathbf{Z}\) -模 \(\mathbf{G}\) ）。映射 \((x, y) \mapsto f(x, y) = g(x \otimes y)\) （从 \(\mathbf{E} \times \mathbf{F}\) 到 \(\mathbf{G}\) ）是 \(\mathbf{Z}\) -双线性的，且满足条件(1)。

\begin{enumerate}
\def\labelenumi{(\alph{enumi})}
\setcounter{enumi}{1}
\tightlist
\item
  反之，设 \(f\) 是一个 \(\mathbf{Z}\) -双线性映射（从 \(\mathbf{E} \times \mathbf{F}\) 到某个 \(\mathbf{Z}\) -模 \(\mathbf{G}\) ），它满足条件(1)。则存在且仅存在一个 \(\mathbf{Z}\) -线性映射 \(g\) （从 \(\mathbf{E} \otimes_{\mathbf{A}} \mathbf{F}\) 到 \(\mathbf{G}\) ），使得 \(f(x,y) = g(x \otimes y)\) 对所有 \(x \in \mathbf{E}, y \in \mathbf{F}\) 成立。
\end{enumerate}

若 \(\phi\) 表示 \(\mathbf{E} \times \mathbf{F}\) 到 \(\mathbf{E} \otimes_{\mathbf{A}} \mathbf{F}\) 的典范映射，则 \(f = g \circ \phi\) ；由此得(a)。为证明(b)，我们注意到，在定义1的记号下， \(f\) 可延拓为一个 \(\mathbf{Z}\) -线性映射 \(\bar{f}\) （从 \(\mathbf{C}\) 到 \(\mathbf{G}\) ）（ \(\S 1\) ，第11号，命题17）。根据关系式(1)， \(\bar{f}\) 在 \(\mathbf{C}\) 的 (2) 中某一类型的所有元素上均为零，从而在 \(\mathbf{D}\) 上为零。因此存在一个 \(\mathbf{Z}\) -线性映射 \(g\) （从 \(\mathbf{C} / \mathbf{D} = \mathbf{E} \otimes_{\mathbf{A}} \mathbf{F}\) 到 \(\mathbf{G}\) ），使得 \(\bar{f} = g \circ \psi\) ，其中 \(\psi: \mathbf{C} \to \mathbf{C} / \mathbf{D}\) 是典范同态（ \(\S 1\) ，第8号，注记）。 \(g\) 的唯一性是直接的，因为 \(\mathbf{E} \otimes_{\mathbf{A}} \mathbf{F}\) 作为 \(\mathbf{Z}\) -模由形如 \(x \otimes y\) 的元素生成。

命题1定义了一个典范同构，它把E × F到G的、满足条件(1)的Z-双线性映射f所成的Z-模映到Z-模Hom \(_{Z}\) (E ⊗A F, G)上。

当 A = Z 时，条件 (1) 对每个 Z-双线性映射 f 自动满足，且 C 的子模 D 已由 (2) 中前两类元素生成。

现在我们回到一般情形，并且 \(E'\) 和 \(F'\) 分别表示E和F的底层Z-模，上述注记和定义1立即表明该Z-模 \(E \otimes_{A} F\) 可以典范地等同于 Z-模 \(E' \otimes_{Z} F'\) 的商模，该商模由形如 \((x\lambda) \otimes y - x \otimes (\lambda y)\) 的元素生成的子 Z-模所决定，其中 x 遍历 E，y 遍历 F，且 \(\lambda\) 遍历A。

推论1. 设H是一个Z-模，且 \(h:E \times F \rightarrow H\) 是一个满足条件(1)的Z-双线性映射，使得H由 \(h(E \times F)\) 生成。假设对于每个Z-模G以及每个从 \(E \times F\) 到G的满足(1)的Z-双线性映射f，

都存在一个Z-线性映射 \(g:H \to G\) ，使得 \(f = g \circ h\) 。那么，若 \(\phi\) 表示 \(E \times F\) 到 \(E \otimes_{A} F\) 的典范映射，则存在且仅存在一个同构 \(\theta\) （从 \(E \otimes_{A} F\) 到H上），使得 \(h = \theta \circ \phi\) 。

假设 \(h(\mathbf{E} \times \mathbf{F})\) 生成 H 蕴含了 g 的唯一性；该推论不过是泛映射问题解的一般唯一性性质（集合论，IV，§ 3，第 1 号）。

推论2. 设 \(\mathbf{E}^0\) （相应地 \(\mathbf{F}^0\) ）表示模 \(\mathbf{E}\) （相应地 \(\mathbf{F}\) ）被视为相反环 \(\mathbf{A}^0\) 上的左（相应地，右）模；则存在且仅存在一个 \(\mathbf{Z}\) -模同构 \(\sigma: \mathbf{E} \otimes_{\mathbf{A}} \mathbf{F} \to \mathbf{F}^0 \otimes_{\mathbf{A}^0} \mathbf{E}^0\) ，使得 \(\sigma(x \otimes y) = y \otimes x\) 对所有 \(x \in \mathbf{E}\) 和 \(y \in \mathbf{F}\) 成立（张量积的“交换性”）。

根据 \(A^{0}\) -模结构（在 \(E^{0}\) 和 \(F^{0}\) 上）的定义，映射 \((x,y)\mapsto y\otimes x\) （从 \(E\times F\) 到 \(F^{0}\otimes_{A^{0}}E^{0}\) ）是Z-双线性的且满足条件(1)，由此得到Z-线性映射 \(\sigma\) 的存在性和唯一性。类似地，可定义一个Z-线性映射 \(\tau:F^{0}\otimes_{A^{0}}E^{0}\to E\otimes_{A}F\) 使得 \(\tau(y\otimes x)=x\otimes y\) ，并且显然 \(\sigma\) 和 \(\tau\) 是互逆的同构。

注记. 非零模的张量积可能为零：例如，取两个Z-模 E = Z/2Z 和 F = Z/3Z，则对所有 \(x \in E\) 和 \(y \in F\) 都有 2x = 0 和 3y = 0；因此，在 E \(\otimes_{\mathbf{Z}} F\) 中，

\[
x \otimes y = 3 (x \otimes y) - 2 (x \otimes y) = x \otimes (3 y) - (2 x) \otimes y = 0
\]

对所有 \(x\) 和 \(y\) 成立（参见第6号，命题6的推论4）。

